% !TEX root = ../main.tex
\clearpage
\phantomsection
\section*{3 \textbf{Literature Review}}
\label{sec:literature-review}
\addcontentsline{toc}{section}{3 Literature Review}

\phantomsection
\subsection*{3.1 Graph-Based On-Chain Reputation Scoring}
\label{sec:graph-based-on-chain-reputation-scoring}
\addcontentsline{toc}{subsection}{3.1 Graph-Based On-Chain Reputation Scoring}

Although prior work often uses Ethereum mainnet as the reference chain, \textbf{this dissertation's primary empirical setting is Arbitrum One with GMX V2 perpetual swap position outcomes} for domain validation. The literature below remains relevant as background for graph-based reputation methods.

Early research on on-chain reputation leveraged graph theory, particularly adaptations of PageRank, to rank wallet addresses by ``trustworthiness.'' Do et al. (2019) provided a first attempt, and subsequent studies formalized this approach for Ethereum. In a recent study, Do et al. (2023) model Ethereum addresses as nodes in a directed graph with transactions as edges, then apply the PageRank algorithm to compute a social credit score for each address. The rationale is that DeFi's growth demands methods to identify and rank entities for unsecured lending and customer protection. In this framework, an address accumulates higher score if it is ``trusted'' by many or large transactions from others, analogous to how important web pages have many incoming links. PageRank thus measures connectivity strength and activity as proxies for reputation. Inspired by the Colendi protocol's vision of a ``social connectivity strength'' metric, these graph models sometimes combine multiple factors (e.g. frequency and volume of interactions) into edge weights. For example, Colendi's system weighted edges by a mix of relation intensity and bilateral activity, feeding into a PageRank-like iterative algorithm to yield a credit score.


A key consideration in graph construction is how to weight the edges. Basic implementations treat each transaction equally, but later work incorporates transaction value as a weight, so that large transfers signify stronger trust. Using a HodgeRank approach, Do et al. (2023) treat the amount of ETH transferred as the weight on each edge and then solve for a global ranking that accounts for these weighted flows. This produces an Ethereum reputation ranking influenced by transaction volume, under the intuition that high-value interactions indicate greater trust or stake between parties. Another practical challenge is dealing with graph structure: Ethereum's transaction graph can contain many isolated clusters or "sink" nodes that absorb PageRank score. Damping factors, as in traditional PageRank, are used to mitigate score inflation for sinks and small isolated subgraphs, ensuring the reputation scores reflect participation in the broader network rather than siloed activity.


One limitation of a vanilla PageRank, i.e., the unmodified PageRank formulation applied to a transaction graph without explicit risk propagation or domain-specific adjustments, is that it may not distinguish malicious but active actors from genuinely reputable ones. Imig (2023) demonstrated that a straightforward PageRank on Ethereum failed to differentiate a known problematic address from a known trustworthy address in terms of score. This highlights the need for refined models. Proposed improvements include incorporating negative feedback or propagation of risk labels through the network. For instance, recent research introduced RiskProp, a network-propagation-based risk scoring model for Ethereum accounts (Yang et al., 2024). RiskProp augments the transaction graph with a "de-anonymous score" for addresses and propagates risk signals through directed bipartite graphs. The approach was able to flag a majority of known high-risk accounts by propagating risk indicators across transaction links, effectively identifying suspicious wallets that a simple connectivity measure might overlook. This illustrates how graph-based methods can be extended beyond pure PageRank to serve as on-chain analogues of credit risk models, where the "credit" in question may be financial reliability or compliance risk.


It is worth noting that graph-based reputation systems often consider temporal dynamics. Real-world creditworthiness changes over time, so some models incorporate time-decay factors to diminish the influence of stale interactions. A dynamic reputation model might weight recent transactions more heavily or periodically expire old edges. This prevents entities from resting on long-past positive interactions and addresses the whitewashing problem, where an entity's bad behavior gets diluted by time. However, striking the right balance is crucial: too much decay can ignore longstanding trust relationships, while too little can allow outdated information to skew the scores.


\phantomsection
\subsection*{3.2 Predictive Modeling for On-Chain Credit Risk}
\label{sec:predictive-modeling-for-on-chain-credit-risk}
\addcontentsline{toc}{subsection}{3.2 Predictive Modeling for On-Chain Credit Risk}

In parallel with graph-theoretical methods, researchers have developed data-driven machine-learning models for on-chain credit scoring, especially in DeFi lending contexts. These models take inspiration from traditional credit analytics, using wallet histories as input features and predicting the likelihood of default or delinquency. Wolf et al. (2022), for example, proposed a credit scoring system for users of the Aave lending platform, built around a supervised machine-learning classifier. In their framework, each wallet's activity is featurized analogously to a credit report: features include payment history, amounts owed, length of credit history, new credit, and credit mix, all derived from on-chain behavior. Using a tree-based ensemble (a binary decision-tree classifier), the model predicts the probability that a given Aave loan position will become delinquent (e.g., fall below collateral requirements and face liquidation) in the near future. The output is essentially a credit risk score for the wallet or loan, indicating the odds of default. Wolf et al. reports that this approach can successfully rank borrowers by risk and thus enable under-collateralized lending by relying on predicted probability of repayment rather than just collateral ratios. Notably, they share an open dataset of Aave user health factors and outcomes to encourage transparency and further research.


Other machine-learning approaches similarly treat on-chain transactions and account data as inputs to predictive models for creditworthiness. Hassija et al. (2020) presented a blockchain-based credit scoring model that leverages prospect theory to weight positive vs. negative credit events differently, capturing the asymmetric risk attitudes of borrowers. Their model, implemented as a smart contract, calculates a trust score for each participant which is updated as new lending outcomes occur.


By incorporating prospect theory (which, for instance, might penalize defaults more heavily than it rewards timely repayments), the scoring mechanism aligns with human risk biases, potentially improving prediction of who will default. Similarly, Jain et al. (2019) developed an Ethereum-based system called Bit-Score for underbanked populations, using a self-sovereign identity to aggregate both financial and non-financial indicators of creditworthiness on-chain. Their design demonstrates the use of alternative data (e.g. social or behavioral information) alongside transaction history to broaden the credit assessment in a decentralized setting.


\phantomsection
\subsection*{3.3 Industry Initiatives and Non-Scholarly Perspectives}
\label{sec:industry-initiatives-and-non-scholarly-perspectives}
\addcontentsline{toc}{subsection}{3.3 Industry Initiatives and Non-Scholarly Perspectives}

Outside the academic literature, the demand for on-chain credit scoring has led to several industry-driven solutions. While not peer-reviewed, these initiatives illustrate real-world interest in the problem and help motivate the EndorseRank research. Cred Protocol is a prominent example: it has developed a "Cred Score" to quantify DeFi lending risk for Ethereum addresses. The Cred Score (range 300-1000, analogous to a FICO credit score) is designed to estimate the probability of loan liquidation or default for a given wallet. Even addresses with no prior borrowing can be scored based on their general on-chain activity patterns. Under the hood, Cred Protocol's model ingests a vast array of features across multiple blockchains and protocols. According to 3Jane Protocol documentation, Cred's scoring engine analyzes 1,000+ features drawn from 54 billion transactions across 8 EVM chains and 100+ DeFi protocols. These include straightforward metrics (loan history, repayment behavior, assets held, debt ratios) as well as more abstract attributes (yield-farming performance, protocol diversity, and participation in governance). A machine-learning model (reportedly trained on over a petabyte of data) synthesizes these inputs into a single risk score. This extensive feature-driven approach aligns with traditional credit bureau analytics, but in a Web3 context. Cred's partners sometimes blend off-chain data as well, and Cred Protocol has integrated with products in the Aave ecosystem to facilitate under-collateralized loans.


Another noteworthy industry example is 3Jane's Jane Score, presented by the protocol as a hybrid on-chain/off-chain credit metric. Jane Score is the native credit measure of the 3Jane protocol, an unsecured DeFi lending platform backed by Paradigm. It consolidates multiple sources: it ingests on-chain credit scores from providers like Cred Protocol and Blockchain Bureau, then combines them with traditional credit bureau scores from TransUnion and Equifax. The result is a composite score (ranging 300-1000) that reflects both the user's blockchain financial history and conventional creditworthiness. For the on-chain component, Jane Score relies on the robust underwriting pipelines described by the protocol, while the off-chain component incorporates factors such as payment history, credit utilization, and length of credit history from VantageScore 3.0 style data. By posting this combined score on-chain with privacy protections via zero-knowledge proofs, 3Jane aims to let smart contracts evaluate a borrower's creditworthiness before issuing an unsecured loan. Although Jane Score and Cred Score originate from industry rather than academia, they underscore the practicality of on-chain credit scoring and inform this research by indicating which features and design choices appear valuable in real deployments.


CredScore and JaneScore also highlight current gaps that the EndorseRank research aims to fill. Notably, existing implementations focus heavily on past transaction behavior (payments, liquidations, balances, etc.), but they do not explicitly utilize the ERC-20 token approval relationships between wallets. An ERC-20 approve (or permit) action, which authorizes another address (typically a smart contract, but sometimes a user's proxy) to spend tokens on one's behalf, is a strong signal of trust or financial linkage that lies outside direct token transfers. For example, if wallet A approves wallet B to spend a large amount of A's tokens, it implies a level of trust or dependence that could influence A's risk: B could empty A's funds if malicious or if B's contract is compromised. Current graph-based scoring systems largely ignore these approval edges, focusing only on realized transfers of value. EndorseRank addresses this research gap by integrating ERC-20 approve/permit edges as first-class links in the reputation graph, treating them as directed weighted edges that can augment a wallet's influence and risk profile in the PageRank computation. In EndorseRank's formulation, if A has approved B for a substantial sum, B might gain ``reputation score'' from A (analogous to a link conferral in PageRank), and A's own score might be influenced by the risk exposure to B. This approach captures implicit credit relationships: approvals often precede financial interactions or indicate delegation of spending power, which are highly relevant to creditworthiness yet missed by transfer-only models.


Furthermore, EndorseRank is distinctive in its handling of temporal effects. Whereas many reputation systems apply time-based decay to down-weight old interactions, EndorseRank proposes to forgo time decay for approval/permit edges. The rationale is that an approval remains in force until revoked; a permission granted a year ago is equally pertinent today if the allowance is still active. By giving persistent weight to long-lived trust links, EndorseRank preserves important information about enduring financial relationships that a short memory window would lose. This contrasts with typical transaction-based scoring, where a payment from years ago might reasonably be given little importance. In the context of approvals, the lack of time decay reflects the enduring nature of the risk: as long as B is authorized to spend A's funds, A's exposure (and hence the reputational linkage) does not diminish over time. This no-decay approach is novel, as most prior works have either not included such edges at all or have not considered the temporal aspect of trust explicitly. By integrating ERC-20 approval networks with PageRank and maintaining their influence over time, EndorseRank is expected to provide a richer and more stable on-chain reputation score---one that captures both the flow of value and the web of delegated permissions on Arbitrum One.


In summary, the literature shows a progression from simple graph-based reputation metrics to sophisticated machine learning models for on-chain credit scoring. PageRank-based techniques contribute a transparent, explainable network-centric view of trust and influence, while ML models offer predictive power for credit risk using diverse features. However, a gap remains in melding the two: incorporating meaningful graph features (like ERC-20 approval links) into a scoring system that can be both explainable and domain-aligned. EndorseRank aims to bridge this gap. It builds on the insight that network structure matters---who trusts whom (via both transfers and approvals)---and that this structure can be quantified through a PageRank-style algorithm. Empirical validation in this dissertation uses GMX V2 perpetual swap position outcomes on Arbitrum One: close-success count and realized gain proxy serve as observable trading-success measures, analogous to the in-degree and in-value proxies used in Do et al.\ (2023). By doing so without arbitrarily discarding older trust links, EndorseRank preserves the long-term context of relationships. This approach directly addresses limitations in prior art and aligns with the real-world needs evidenced by industry solutions like Cred Score and Jane Score. It promises an enhanced on-chain wallet reputation score that could improve risk assessment for leveraged DeFi activity, making decentralized finance both more accessible and more secure.
